\documentclass[11pt,a4paper]{article}
\usepackage[utf8]{inputenc}
\usepackage[T1]{fontenc}
\usepackage[swedish]{babel}
\usepackage{amsmath,amssymb}
\usepackage{geometry}
\usepackage{fancyhdr}
\usepackage{array}
\usepackage{booktabs}
\usepackage{xcolor}
\usepackage{tcolorbox}
\usepackage{hyperref}

\geometry{margin=2.5cm}
\pagestyle{fancy}
\fancyhf{}
\fancyhead[L]{Hvitfeldska Spetsutbildning}
\fancyhead[R]{Matematikhandbok 2026}
\fancyfoot[C]{\thepage}

\definecolor{mainblue}{RGB}{26,84,144}
\definecolor{lightblue}{RGB}{240,244,248}

\tcbset{
    colback=lightblue,
    colframe=mainblue,
    boxrule=1pt,
    arc=3mm
}

\title{\Huge\textbf{HVITFELDTSKAS SPETSUTBILDNING}\\[0.5cm]
\Large Komplett Matematikhandbok 2026}
\author{Skapad for Yvonna}
\date{Provdatum: 24 februari 2026}

\begin{document}

\maketitle
\tableofcontents
\newpage

\section{TALTEORI}

\subsection{Delbarhetsregler}

\begin{tabular}{|c|l|l|}
\hline
\textbf{Tal} & \textbf{Regel} & \textbf{Exempel} \\
\hline
2 & Sista siffran ar jamn & 1234 \\
3 & Siffersumman delbar med 3 & 156: 1+5+6=12 \\
4 & Sista tva siffrorna delbara med 4 & 1324: 24/4=6 \\
5 & Slutar pa 0 eller 5 & 1235 \\
6 & Delbar med bade 2 och 3 & 1254 \\
8 & Sista tre siffrorna delbara med 8 & 5016 \\
9 & Siffersumman delbar med 9 & 2286: 2+2+8+6=18 \\
10 & Slutar pa 0 & 1230 \\
\hline
\end{tabular}

\subsection{Primtal}

\textbf{Definition:} Ett primtal ar ett heltal $> 1$ som endast ar delbart med 1 och sig sjalv.

\begin{tcolorbox}[title=Primtal under 100]
2, 3, 5, 7, 11, 13, 17, 19, 23, 29, 31, 37, 41, 43, 47, 53, 59, 61, 67, 71, 73, 79, 83, 89, 97
\end{tcolorbox}

\subsection{SGD och MGM}

\textbf{SGD} (Storsta Gemensamma Delare): Ta gemensamma primfaktorer med \textbf{lagsta} exponent.

\textbf{MGM} (Minsta Gemensamma Multipel): Ta alla primfaktorer med \textbf{hogsta} exponent.

\begin{tcolorbox}[title=Viktig formel]
\[ \text{SGD}(a,b) \times \text{MGM}(a,b) = a \times b \]
\end{tcolorbox}

\subsection{TYPUPPGIFT: Heltalsekvationer}

\textbf{Problem:} $x, y, z$ ar heltal $> 1$. $xy = 350$, $yz = 180$, $xz$ ar delbar med 8. Bestam $xz$.

\textbf{Losning:}
\begin{enumerate}
\item Primtalsfaktorisera: $350 = 2 \times 5^2 \times 7$, $180 = 2^2 \times 3^2 \times 5$
\item SGD$(350, 180) = 2 \times 5 = 10$. Mojliga $y$: 2, 5, 10
\item Testa $y = 5$: $x = 70$, $z = 36$, $xz = 2520$
\item Kontroll: $2520 / 8 = 315$ (heltal!)
\end{enumerate}

\begin{tcolorbox}[colback=green!10, colframe=green!50!black]
\textbf{SVAR:} $xz = 2520$
\end{tcolorbox}

\section{ALGEBRA}

\subsection{Konjugat- och Kvadreringsregler}

\begin{align}
a^2 - b^2 &= (a+b)(a-b) \\
(a+b)^2 &= a^2 + 2ab + b^2 \\
(a-b)^2 &= a^2 - 2ab + b^2
\end{align}

\subsection{Potensregler}

\begin{tabular}{|l|l|}
\hline
\textbf{Regel} & \textbf{Formel} \\
\hline
Multiplikation & $a^m \cdot a^n = a^{m+n}$ \\
Division & $a^m / a^n = a^{m-n}$ \\
Potens av potens & $(a^m)^n = a^{mn}$ \\
Negativ exponent & $a^{-n} = 1/a^n$ \\
Nollexponent & $a^0 = 1$ \\
\hline
\end{tabular}

\subsection{Andragradsekvationer}

\begin{tcolorbox}[title=PQ-formeln]
For $x^2 + px + q = 0$:
\[ x = -\frac{p}{2} \pm \sqrt{\left(\frac{p}{2}\right)^2 - q} \]
\end{tcolorbox}

\begin{tcolorbox}[title=ABC-formeln]
For $ax^2 + bx + c = 0$:
\[ x = \frac{-b \pm \sqrt{b^2 - 4ac}}{2a} \]
\end{tcolorbox}

\subsection{Talföljder}

\textbf{Aritmetisk talföljd:}
\begin{align}
a_n &= a_1 + (n-1)d \\
S_n &= \frac{n(a_1 + a_n)}{2}
\end{align}

\textbf{Geometrisk talföljd:}
\begin{align}
a_n &= a_1 \cdot k^{n-1} \\
S_n &= a_1 \cdot \frac{k^n - 1}{k - 1}
\end{align}

\section{GEOMETRI}

\subsection{Trianglar}

\textbf{Vinkelsumma:} Summan av vinklarna = $180°$

\begin{tcolorbox}[title=Pythagoras sats]
I en ratvinklig triangel med kateter $a$, $b$ och hypotenusa $c$:
\[ a^2 + b^2 = c^2 \]
\end{tcolorbox}

\textbf{Pythagoriska tripplar:} 3-4-5, 5-12-13, 8-15-17, 7-24-25

\textbf{Triangelns area:}
\[ A = \frac{bh}{2} \]

\subsection{Speciella trianglar}

\textbf{Liksidig triangel} (sida $a$):
\begin{itemize}
\item Hojd: $h = \frac{a\sqrt{3}}{2}$
\item Area: $A = \frac{a^2\sqrt{3}}{4}$
\end{itemize}

\textbf{30-60-90 triangel:} Sidor i forhallande $1 : \sqrt{3} : 2$

\textbf{45-45-90 triangel:} Sidor i forhallande $1 : 1 : \sqrt{2}$

\subsection{Cirkeln}

\begin{tabular}{|l|l|}
\hline
\textbf{Storhet} & \textbf{Formel} \\
\hline
Omkrets & $O = 2\pi r = \pi d$ \\
Area & $A = \pi r^2$ \\
Baglangd & $b = \frac{v}{360} \cdot 2\pi r$ \\
Sektorarea & $A = \frac{v}{360} \cdot \pi r^2$ \\
\hline
\end{tabular}

\begin{tcolorbox}[colback=yellow!10, colframe=orange]
\textbf{Viktigt:} Tangenten ar vinkelrat mot radien vid beroringspunkten!
\end{tcolorbox}

\subsection{Koordinatgeometri}

\textbf{Avstand mellan tva punkter:}
\[ d = \sqrt{(x_2-x_1)^2 + (y_2-y_1)^2} \]

\textbf{Mittpunkt:}
\[ M = \left(\frac{x_1+x_2}{2}, \frac{y_1+y_2}{2}\right) \]

\textbf{Lutning:}
\[ k = \frac{y_2 - y_1}{x_2 - x_1} \]

\textbf{Vinkelrata linjer:} $k_1 \cdot k_2 = -1$

\section{KOMBINATORIK}

\subsection{Grundprinciper}

\textbf{Multiplikationsprincipen:} $n_1 \times n_2 \times \ldots \times n_k$ satt

\textbf{Additionsprincipen:} $n + m$ satt (om ej samtidiga)

\textbf{Komplementprincipen:} Gynnsamma = Totalt $-$ Ogynnsamma

\subsection{Permutationer och Kombinationer}

\textbf{Fakultet:}
\[ n! = n \times (n-1) \times \ldots \times 2 \times 1 \]

\textbf{Permutation} (ordnat urval):
\[ P(n,r) = \frac{n!}{(n-r)!} \]

\textbf{Kombination} (oordnat urval):
\[ C(n,r) = \binom{n}{r} = \frac{n!}{r!(n-r)!} \]

\subsection{TYPUPPGIFT: Vagrakning i rutnat}

\textbf{Problem:} Ga fran S till M i rutnat. Endast uppat eller hoger.

\begin{tcolorbox}[title=Losning]
Om du gar $m$ steg hoger och $n$ steg uppat:
\[ \text{Antal vagar} = \binom{m+n}{m} = \binom{m+n}{n} \]
\end{tcolorbox}

\textbf{Undvik punkt P:} Antal = Totalt $-$ (Vagar via P)

\section{SANNOLIKHET}

\begin{tcolorbox}[title=Grundformel]
\[ P(A) = \frac{\text{Antal gynnsamma utfall}}{\text{Antal mojliga utfall}} \]
\end{tcolorbox}

\textbf{Rakneregler:}
\begin{itemize}
\item Komplement: $P(A') = 1 - P(A)$
\item Addition: $P(A \cup B) = P(A) + P(B) - P(A \cap B)$
\item Multiplikation (oberoende): $P(A \cap B) = P(A) \times P(B)$
\end{itemize}

\begin{tcolorbox}[colback=green!10, colframe=green!50!black]
\textbf{Minnesregel:} P(minst en) = 1 $-$ P(ingen)
\end{tcolorbox}

\section{ORDPROBLEM}

\subsection{Samarbetsproblem}

\begin{tcolorbox}[title=Formel: Tva arbetare]
Om A gor arbetet pa $a$ timmar och B pa $b$ timmar:
\[ t = \frac{ab}{a+b} \]
\end{tcolorbox}

\subsection{Rorelseproblem}

\textbf{Grundformel:} $s = v \cdot t$

\textbf{Motas:} $t = \frac{\text{stracka}}{v_1 + v_2}$

\textbf{Ikapp:} $t = \frac{\text{forsprang}}{v_{\text{snabb}} - v_{\text{langsam}}}$

\section{FORMELSAMLING}

\subsection{Algebra}
\begin{align*}
(a+b)^2 &= a^2 + 2ab + b^2 \\
(a-b)^2 &= a^2 - 2ab + b^2 \\
a^2 - b^2 &= (a+b)(a-b)
\end{align*}

\subsection{Geometri}
\begin{align*}
\text{Pythagoras:} \quad & a^2 + b^2 = c^2 \\
\text{Triangelarea:} \quad & A = \frac{bh}{2} \\
\text{Cirkelomkrets:} \quad & O = 2\pi r \\
\text{Cirkelarea:} \quad & A = \pi r^2
\end{align*}

\subsection{Kombinatorik}
\begin{align*}
\text{Kombination:} \quad & \binom{n}{r} = \frac{n!}{r!(n-r)!} \\
\text{Vagrakning:} \quad & \binom{m+n}{m}
\end{align*}

\section{MINNESREGLER}

\begin{itemize}
\item \textbf{Talteori:} "3 och 9: Siffersumma!"
\item \textbf{SGD/MGM:} "SGD = Minsta exponenter, MGM = Storsta"
\item \textbf{Kombinatorik:} "Ordning spelar roll? Permutation. Annars Kombination."
\item \textbf{Vagrakning:} "Hoger + Uppat = Kombination!"
\item \textbf{Geometri:} "Tangent vinkelrat mot radie"
\end{itemize}

\vfill
\begin{center}
\textbf{\Large Lycka till pa provet, Yvonna!}

Version 1.0 -- December 2024
\end{center}

\end{document}
